\chapter{Generalization, Regularization, and Training}

A security team builds an intrusion detector from two raw quantities: how many
requests per second and the average payload size. From those two quantities they
derive twenty-one features by multiplying and raising both to the fifth power,
then train a logistic classifier on sixty labeled connections. On the training
data, the detector is wrong on zero of the sixty connections. On twenty thousand
new connections it has never seen, it is wrong on thirty-two percent. The same
model, the same weights, two numbers far apart.

That zero is neither a coincidence nor an achievement. With twenty-one features
and sixty examples, a separating plane can always be found that cuts every
training example correctly, including the example whose label is wrong. What the
model learns is not the shape of the correct separator but the position of those
sixty points. As soon as the points change, the plane is in the wrong place.

Chapter 25 taught us to compute gradients and follow them down until the training
error is small; Chapter 20 taught us to read a single number together with its
uncertainty. Neither answers the question that decides matters here: which
quantity must be made small so that the number in the field shrinks as well, and
how to choose training settings without deceiving ourselves.

\begin{untukanak}
A student is given sixty practice problems before an exam. He studies by
memorizing: every problem together with its answer, down to the letter. On the
practice problems he is always right, a hundred out of a hundred, and he goes
home proud. On the exam the problems are different: different numbers, different
sentences, different order. His grade collapses, and the teacher is puzzled,
because this student looked like the most diligent one.

A second student learns more slowly. He does not memorize answers; he looks for
the pattern: when to divide, when to multiply, when an answer makes sense. On the
practice problems he gets several wrong, and he is not proud. On the exam he does
better.

Now the part that makes one curious. Suppose the first student were given two
thousand practice problems instead of sixty. Is memorizing still bad? It turns
out that it is not always: if there are enough problems and they vary enough,
memorizing becomes an expensive way of learning that works. So the problem is not
memorizing but memorizing from few examples. In this chapter we will learn why
the number of examples and the complexity of the learning method must be
balanced, why adding complexity actually makes the exam grade worse, and how to
measure that balance without waiting for the exam.
\end{untukanak}

\section{Two errors and one gap}
\tingkat{basic}

Every trained model has two numbers, and both are computed in exactly the same
way: the proportion of mistakes on a set of examples. The only difference is
which set is used. In everyday language: the grade on a test at home differs from
the grade on a test at school, even though the problems are of the same kind.

\kamus{generalization}{generalisasi}{the ability to work on new problems, not on problems
that have already been memorized}

\begin{definisi}[Training error, population error, and the gap]
\index{generalization}
Suppose the examples $(x_i,y_i)$ are drawn from an unknown distribution $\mathcal{D}$.
For a hypothesis $h$ and 0-1 error, the \emph{population error} is
$R(h)=\mathbb{P}_{(x,y)\sim\mathcal{D}}\bigl(h(x)\neq y\bigr)$, whereas the
\emph{training error} on $n$ examples is
$\hat R_n(h)=\frac1n\sum_{i=1}^n \mathbf{1}\{h(x_i)\neq y_i\}$.
The difference $R(h)-\hat R_n(h)$ is called the \emph{generalization gap} for that $h$.
\end{definisi}

Training error can be computed; population error cannot, because the distribution
is unknown. All we can do is estimate it on another set that is not used for
training, and that estimate itself wobbles like every proportion in Chapter 20. So
every "validation error" number in this chapter is an estimate, not the truth.

Now two opposite diseases. In everyday language: a model that is too simple cannot
capture the pattern, whereas a model that is too complex memorizes coincidence.

\begin{definisi}[Underfitting and overfitting]
\index{underfitting}
\index{overfitting}
A model \emph{underfits} when its training error and its population error are both
high: the model is too rigid for the actual shape. A model \emph{overfits} when its
training error is low but its generalization gap is large: the model adapts itself
to the coincidence in the training data, not to the actual shape.
\end{definisi}

The two diseases are not symmetric in practice. Underfitting is visible: the
training error is high, and adding features helps. Overfitting deceives, because
the only number that is easily available, namely the training error, is precisely
the one that looks perfect. The model in the chapter opening is the purest
example: its largest weight is $406.204$, a sign that its score depends on small
differences between points, not on a plausible scale.

What governs both is \emph{capacity}: how rich the set of functions is that the
model may choose from. For a linear classifier over $d$ features, capacity rises
with $d$; for a neural network, with the number of units, the depth, and the number
of training steps. Capacity is not merely a parameter count: a large model with
strong regularization can behave like a small model.

\begin{bacakeras}
Two voices, alternating.

``A hundred out of a hundred. That means he understands, doesn't it?''

``Not necessarily. A hundred out of a hundred only means there was no new problem.''

``So memorizing is wrong?''

``Memorizing is not wrong. Memorizing from few examples is.''

``Where is the difference?''

``In the problems that are not in the notes. Memorization answers what has already
been seen. Understanding answers what has not been seen.''

``And if there are two thousand problems instead of sixty?''

``Then memorization starts to help. Not because memorization turns into
understanding, but because the stored coincidences are used again more often.''

Say it once more, slowly: memorization answers what is already there, understanding
answers what is not yet there. Repeat that line until it feels natural in the
mouth: already there, not yet there.
\end{bacakeras}

The generalization gap shrinks as examples are added. In the capacity experiment
with a fixed step budget, the gap is $0.1338$ at $n=60$, $0.0274$ at $n=150$,
$0.0193$ at $n=600$, and $0.0097$ at $n=2400$, whereas the training error barely
moves, staying in the range $0.17$ to $0.21$.

\section{The bias-variance decomposition}
\tingkat{intermediate}

A sharper question: where the gap comes from. In everyday language: a shooter can
miss because the sights are shifted to one side, or because his hands shake, and
the two causes call for different remedies.

\kamus{bias-variance}{bias-variansi}{a systematic miss versus shaking; both contribute
error, and pressing one down sometimes raises the other}

\begin{teorema}[Bias-variance decomposition]
\index{bias-variance}
Suppose $y=f(x)+\varepsilon$ with $\mathbb{E}[\varepsilon]=0$ and
$\operatorname{Var}(\varepsilon)=\sigma^2$, and suppose $\hat f$ is trained on a
random training set $D$ that is independent of $\varepsilon$. For every fixed $x$,
\[
\mathbb{E}_{D,\varepsilon}\bigl[(y-\hat f(x))^2\bigr]
=\sigma^2
+\bigl(\mathbb{E}_D\hat f(x)-f(x)\bigr)^2
+\operatorname{Var}_D\bigl(\hat f(x)\bigr).
\]
\end{teorema}

\begin{proof}[Proof]
Write $y-\hat f(x)=\varepsilon+\bigl(f(x)-\mathbb{E}_D\hat f(x)\bigr)
+\bigl(\mathbb{E}_D\hat f(x)-\hat f(x)\bigr)$, namely three terms: noise that cannot
be reduced, bias, and the deviation from the mean. Square it and take the
expectation. The cross terms containing $\varepsilon$ vanish because
$\mathbb{E}[\varepsilon]=0$ and $\varepsilon$ is independent of $D$; the cross term
between bias and deviation vanishes because
$\mathbb{E}_D[\hat f(x)-\mathbb{E}_D\hat f(x)]=0$. What remains is exactly the three
terms of the theorem's statement.
\end{proof}

The first term cannot be reduced by any model; it belongs to the data, is called
the irreducible error, and on the data of the chapter opening its value is about
$0.21$ because the labels are random. The second and third terms can be traded:
adding capacity lowers bias and raises variance, reducing capacity does the
opposite. Regularization is the name for that trade.

The cheapest example that contains both terms at once is the shrinkage estimator.
To measure $\mu$ from $n$ noisy observations with variance $\sigma^2$, people
usually use $\bar X_n$, which is unbiased. We may use $c\bar X_n$ for $c\in(0,1)$,
and that estimator is deliberately biased. Its mean squared error
\[
\operatorname{MSE}(c)=(c-1)^2\mu^2+c^2\frac{\sigma^2}{n},
\]
is minimized at $c^\ast=\mu^2/(\mu^2+\sigma^2/n)$. For $\mu=2$, $\sigma=1$, and
$n=25$ we obtain $c^\ast=0.990099$: the best estimator is in fact slightly shrunk.

\begin{center}
\small
\begin{tabular}{rrrr}
\hline
$c$ & bias$^2$ & variance & MSE \\
\hline
0.70 & 0.360000 & 0.019600 & 0.379600 \\
0.90 & 0.040000 & 0.032400 & 0.072400 \\
0.99 & 0.000400 & 0.039204 & 0.039604 \\
1.00 & 0.000000 & 0.040000 & 0.040000 \\
1.10 & 0.040000 & 0.048400 & 0.088400 \\
1.30 & 0.360000 & 0.067600 & 0.427600 \\
\hline
\end{tabular}
\end{center}

Three things must be read from that table. The gain from shrinking is small: MSE
falls from $0.040000$ to $0.039604$, about one percent. The loss from deviating in
the wrong direction is far larger: $c=1.3$ gives $0.4276$, whereas $c=0.7$ gives
only $0.3796$. And the minimum is flat, so choosing a slightly wrong $c$ is not a
disaster.

The theorem does not give a formula for an $\hat f$ as complicated as a neural
network. Its decomposition holds formally, but neither term can be computed from
data alone; what is observed is only their sum, namely the validation error. For
that reason the rest of this chapter works with observed numbers, not with the true
bias and variance.
\section{L2 and L1 regularization}
\tingkat{intermediate}

The most direct way to press down variance is to forbid the weights from growing
large. In everyday language: we add a fine to the objective, so that the model has
to pay for every large weight, and will pay only if the gain on the training data
is worth it.

\kamus{regularization}{regularisasi}{a fine that the model pays for large weights, so that
it does not memorize the training data}

\begin{definisi}[The regularized objective]
\index{regularization}
For an original objective $J(\mathbf{w})$ and a penalty $\Omega$, the objective being
minimized is $\tilde J(\mathbf{w})=J(\mathbf{w})+\lambda\Omega(\mathbf{w})$ with
$\lambda\ge0$ as a hyperparameter. For the L2 penalty, $\Omega=\frac12\|\mathbf{w}\|_2^2$;
for the L1 penalty, $\Omega=\|\mathbf{w}\|_1=\sum_j|w_j|$.
\end{definisi}

For L2, the gradient of the objective is $\nabla J+\lambda\mathbf{w}$, so one
gradient-descent step with rate $\eta$ has the form
\[
\mathbf{w}\leftarrow\mathbf{w}-\eta\bigl(\nabla J+\lambda\mathbf{w}\bigr)
=(1-\eta\lambda)\mathbf{w}-\eta\nabla J .
\]
The factor $(1-\eta\lambda)$ is the origin of the name \emph{weight decay}: at every
step, before the gradient is used, the weights are multiplied by a number slightly
less than one. If the gradient is zero, the weights decay toward zero geometrically
with half-life $\ln 2/(\eta\lambda)$ steps.

\begin{akibat}[Weight decay is the same as an L2 penalty]
\index{weight decay}
For a rate $\eta$ and an L2 penalty with coefficient $\lambda$, gradient descent on
the regularized objective is identical to gradient descent on the original
objective followed by multiplying the weights by $1-\eta\lambda$. Thus $\lambda$ and
$\eta$ appear only as the product $\eta\lambda$ in the decay step, so doubling the
rate and halving $\lambda$ give the same decay trajectory.
\end{akibat}

The L2 penalty also has a Bayesian reading: it is equivalent to a Gaussian prior on
the weights, so the solution is a MAP estimate, not a maximum-likelihood one. That
explains why L2 believes small weights are more likely, and why it never forces a
single weight to zero.

L1 behaves differently, and the whole difference lies in the shape of its penalty
near zero. In everyday language: a quadratic penalty pulls ever more weakly once the
weight is small, whereas an absolute-value penalty pulls with constant strength
until the weight is exactly zero.

\begin{lemma}[Soft thresholding]
\index{regularization}
For $a\in\mathbb{R}$ and $t>0$, the minimizer of $u\mapsto\frac12(u-a)^2+t|u|$ is
\[
u^\ast=\operatorname{sign}(a)\max\{|a|-t,0\},
\]
so that $u^\ast=0$ exactly when $|a|\le t$.
\end{lemma}

For $u>0$ the function is $\frac12(u-a)^2+tu$ with derivative $u-a+t$, so its
critical point is $u=a-t$, valid only when $a>t$. For $u<0$ the derivative is
$u-a-t$ with critical point $u=a+t$, valid only when $a<-t$. At $u=0$ the
subgradient is the interval $[-a-t,-a+t]$, which contains zero exactly when
$|a|\le t$; in that case zero is the minimizer because the function is convex. The
three cases join into the formula above.

That lemma explains the whole behavior of L1. With a quadratic approximation around
the unpenalized solution $w^\ast$ and curvature $H$, the objective is approximately
$\frac H2(w-w^\ast)^2+\lambda|w|$, so that
$w=\operatorname{sign}(w^\ast)\max\{|w^\ast|-\lambda/H,0\}$: every weight below
$\lambda/H$ becomes exactly zero, not merely small. For L2 the solution is
$w=w^\ast/(1+\lambda/H)$, which is never zero unless $w^\ast$ itself is zero.

\begin{proposisi}[Why L1 produces true zeros]
\index{regularization}
Suppose $J$ is differentiable and convex around zero. The point $\mathbf{w}=\mathbf{0}$
is a minimizer of $J(\mathbf{w})+\lambda\|\mathbf{w}\|_1$ if and only if
$|\partial J/\partial w_j|\le\lambda$ for every $j$. For the L2 penalty, the same
condition is $\partial J/\partial w_j=0$ for every $j$.
\end{proposisi}

\begin{proof}[Proof]
The subgradient of $\|\cdot\|_1$ at zero is the box $[-1,1]^d$, so the condition
$0\in\nabla J(\mathbf{0})+\lambda[-1,1]^d$ is equivalent to $|(\nabla J)_j|\le\lambda$
for all $j$; for a convex function, that condition is necessary and sufficient. For
the L2 penalty, the gradient is $\nabla J+\lambda\mathbf{w}$, which at $\mathbf{0}$
equals $\nabla J$, so the condition is $\nabla J(\mathbf{0})=\mathbf{0}$.
\end{proof}

The practical consequence matters for reading results. L1 gives a flat region of
width $2\lambda$ around zero, so that weights that are "almost useless" are switched
off entirely; L2 only shrinks them. For that reason L1 is used as a feature
selector, and for the same reason the number of zero weights under L1 is not a
delicate measure of importance: it jumps, it does not shift slowly.

Now the numbers. The same logistic classifier, the same data, twenty-one features
on sixty training examples, everything computed with the script in the lab section.

\begin{center}
\small
\begin{tabular}{rrrr}
\hline
$\lambda$ & $\max_j|w_j|$ & training error & validation error \\
\hline
0       & 406.204 & 0.0000 & 0.3262 \\
0.001   &   3.056 & 0.1000 & 0.2927 \\
0.003   &   1.814 & 0.1000 & 0.2688 \\
0.010   &   1.018 & 0.1167 & 0.2475 \\
0.030   &   0.631 & 0.1167 & 0.2433 \\
0.100   &   0.399 & 0.1500 & 0.2406 \\
0.300   &   0.268 & 0.2333 & 0.2452 \\
1.000   &   0.142 & 0.2500 & 0.2695 \\
3.000   &   0.062 & 0.3333 & 0.3205 \\
10.000  &   0.022 & 0.3333 & 0.3395 \\
100.000 &   0.003 & 0.3333 & 0.3395 \\
\hline
\end{tabular}
\end{center}

\begin{figure}[htbp]
\centering
\begin{tikzpicture}
\begin{axis}[
  width=12.2cm, height=5.6cm,
  xmin=-0.4, xmax=10.4, ymin=0, ymax=0.40,
  axis lines=left,
  xtick={0,3,5,7,10},
  xticklabels={$0$,$10^{-2}$,$0.1$,$1$,$100$},
  ytick={0,0.1,0.2,0.3,0.4},
  yticklabels={$0$,$0.10$,$0.20$,$0.30$,$0.40$},
  tick label style={font=\scriptsize},
  xlabel={regularization strength $\lambda$ (capacity decreases to the right)},
  ylabel={error},
  xlabel style={font=\small},
  ylabel style={font=\small},
  legend style={font=\scriptsize,draw=none,at={(0.03,0.97)},anchor=north west},
]
\addplot[very thick,layerindigo,mark=*,mark size=1.5pt] coordinates {
  (0,0.0000) (1,0.1000) (2,0.1000) (3,0.1167) (4,0.1167) (5,0.1500)
  (6,0.2333) (7,0.2500) (8,0.3333) (9,0.3333) (10,0.3333)};
\addlegendentry{training error}
\addplot[very thick,accent,mark=square*,mark size=1.5pt] coordinates {
  (0,0.3262) (1,0.2927) (2,0.2688) (3,0.2475) (4,0.2433) (5,0.2406)
  (6,0.2452) (7,0.2695) (8,0.3205) (9,0.3395) (10,0.3395)};
\addlegendentry{validation error}
\addplot[only marks,mark=o,mark size=3.4pt,accent] coordinates {(5,0.2406)};
\draw[densely dashed,layergrey] (axis cs:5,0)--(axis cs:5,0.2406);
\node[accent,font=\scriptsize,anchor=north] at (axis cs:4.55,0.238)
  {minimum $0.2406$};
\end{axis}
\end{tikzpicture}
\caption{Training error and validation error against the regularization strength
$\lambda$; the direction to the right means capacity decreases, because the weights
are held back ever more tightly. The eleven points are the rows of the previous
table. The training error rises monotonically from $0.0000$ to $0.3333$, whereas the
validation error forms a valley: $0.3262$ at $\lambda=0$, a minimum of $0.2406$ at
$\lambda=0.1$, then $0.3395$ at $\lambda=10$ and above. The bottom of the valley is
wide, from $0.2433$ at $\lambda=0.03$ to $0.2452$ at $\lambda=0.3$.}
\label{fig:lembah-kapasitas}
\end{figure}

Figure~\ref{fig:lembah-kapasitas} plots the third and fourth columns of that table
as a single valley. Read that table as three regions, not eleven rows. Without
regularization the training error is zero and the validation error is $0.3262$, with
weights up to the hundreds because nothing holds them back. In the region of
$\lambda$ between $0.001$ and $0.1$, the training error rises to $0.1500$ while the
validation error falls to $0.2406$. In the region $\lambda\ge3$ both are bad:
$0.3333$ and $0.3395$.

The last row deserves emphasis. With $\lambda=10$ the validation error is $0.3395$,
\emph{worse} than the $0.3262$ of the model without regularization. So regularization
is not a medicine that always helps; it helps at a certain dose and does damage
beyond it. At $\lambda=10$ the model practically always guesses the majority class,
whose error is $0.3419$: its weights are so small that the score is almost entirely
determined by the bias.

The L1 penalty gives a similar picture with a different shape. At $\lambda=0.05$ its
validation error is $0.2334$ with eleven of the twenty-one weights alive; at
$\lambda=0.2$ only one weight remains, with error $0.3395$. The count of live weights
does not fall monotonically: $\lambda=0.005$ leaves eleven, $\lambda=0.01$ leaves
fourteen. That number is stable against a fourfold increase in iterations, so it is
not an optimizer failure; the L1 trajectory for the logistic loss really does not
guarantee a monotone count of zeros, and that is one more reason not to read it as a
measure of importance.

The final comparison: L1 reaches $0.2334$ and L2 $0.2406$. The difference $0.0072$ is
smaller than the wobble between partitions, about $0.02$, so this data is not enough
to conclude that L1 is better; all that can be concluded is that both improve on the
model without regularization.
\begin{untukanak}
Imagine a class that has to guess a friend's body weight. The teacher gives two
methods. Method one: every guess that is too far from the answer is fined an amount
that grows the farther away it is. Method two: every guess that is not exactly the
answer is fined a fixed amount, no matter whether the difference is small or large.

Method one makes all guesses shift slightly toward the middle, but none of them
becomes zero. Method two rounds guesses that are already close straight down to
zero, because the fine is fixed and no longer worth the improvement that might be
gained. That is the difference between two fines that look alike: one shrinks, the
other switches off.

The question: if we have twenty-one features and only sixty examples, what kind of
fine makes the model stop memorizing? The answer is not "a large fine", because a
fine that is too large switches everything off, including what is correct. In this
chapter we will learn why there is a dose, why that dose must be searched for, and
why a dose that is searched for on the wrong data will deceive us neatly.
\end{untukanak}

\section{Early stopping, dropout, and augmentation}
\tingkat{intermediate}

Three other ways of pressing down variance without touching the formula of the
objective. In everyday language: we can stop learning before coincidence begins to
be memorized, we can switch off part of the memory at random, and we can multiply
examples by changing the examples we already have.

Early stopping is the cheapest form of regularization, because it adds not a single
line to the objective: it only stops the trajectory. During training, the training
error almost always falls, but the validation error falls and then rises. The lowest
point of the validation error is the one taken, and the number of steps becomes a
hyperparameter to be searched for.

\begin{center}
\small
\begin{tabular}{rrrr}
\hline
epoch & training error & validation error & $\max_j|w_j|$ \\
\hline
1    & 0.1333 & 0.2481 & 0.15 \\
5    & 0.1333 & 0.2429 & 0.38 \\
10   & 0.1333 & 0.2432 & 0.48 \\
50   & 0.1167 & 0.2454 & 0.77 \\
200  & 0.1167 & 0.2541 & 1.38 \\
800  & 0.1000 & 0.2879 & 2.85 \\
2400 & 0.0833 & 0.3109 & 4.81 \\
\hline
\end{tabular}
\end{center}

\begin{figure}[htbp]
\centering
\begin{tikzpicture}
\begin{axis}[
  width=11.8cm, height=5.6cm,
  xmode=log,
  xmin=1, xmax=3000, ymin=0.05, ymax=0.35,
  axis lines=left,
  xtick={1,10,100,1000},
  xticklabels={$1$,$10$,$100$,$1000$},
  ytick={0.10,0.15,0.20,0.25,0.30},
  yticklabels={$0.10$,$0.15$,$0.20$,$0.25$,$0.30$},
  tick label style={font=\scriptsize},
  xlabel={training epoch (log scale)},
  ylabel={error},
  xlabel style={font=\small},
  ylabel style={font=\small},
  legend style={font=\scriptsize,draw=none,at={(0.03,0.97)},anchor=north west},
]
\addplot[very thick,layerteal,mark=*,mark size=1.5pt] coordinates {
  (1,0.1333) (5,0.1333) (10,0.1333) (50,0.1167) (200,0.1167) (800,0.1000)
  (2400,0.0833)};
\addlegendentry{training error}
\addplot[very thick,accent,mark=square*,mark size=1.5pt] coordinates {
  (1,0.2481) (5,0.2429) (10,0.2432) (50,0.2454) (200,0.2541) (800,0.2879)
  (2400,0.3109)};
\addlegendentry{validation error}
\draw[densely dashed,layergrey] (axis cs:8,0.05)--(axis cs:8,0.2427);
\addplot[only marks,mark=o,mark size=3.4pt,accent] coordinates {(8,0.2427)};
\node[accent,font=\scriptsize,anchor=north west] at (axis cs:8.6,0.2325)
  {epoch 8: $0.2427$};
\end{axis}
\end{tikzpicture}
\caption{Both errors along one training trajectory, horizontal axis on a log scale.
The training error falls without ever rising, from $0.1333$ at epoch 1 to $0.0833$
at epoch 2400. The validation error falls to $0.2427$ at epoch 8 and then rises to
$0.3109$. The dashed line marks the best stopping point: stopping there gives an
error $0.0682$ lower than training all the way to epoch 2400, and the curves begin to
separate just after that point.}
\label{fig:dua-kurva}
\end{figure}

Figure~\ref{fig:dua-kurva} draws those two columns. The training error falls along
the trajectory, from $0.1333$ to $0.0833$, without ever rising, whereas the
validation error falls until epoch 8 with a best value of $0.2427$ and then rises to
$0.3109$ at epoch 2400. The last column explains why: what grows throughout training
is the size of the weights, and that is the variance we are paying for.

Early stopping is not merely a tactic. For a quadratic objective and a small rate,
stopping at step $\tau$ is nearly equivalent to an L2 penalty with coefficient about
$1/(\tau\eta)$; that is an approximation, not an exact identity. For that reason the
number of steps is a capacity hyperparameter, not a matter of "training until
finished".

\begin{resep}
\begin{enumerate}
  \item Split three sets before a single training step is run: train, validation,
        test. The test set is sealed.
  \item Train on the training set only; every few epochs compute the error on the
        validation set, and record the three values (epoch, validation error,
        training error).
  \item Take the epoch with the smallest validation error, then return the weights to
        that point; do not use the weights of the last epoch.
  \item Only now open the test set, compute its error once, and report it together
        with the chosen epoch and the width of the validation curve.
\end{enumerate}
A warning about order: step 1 must be finished before step 2, and step 4 must not be
used to shift the choice made in step 3. If the test set is looked at while the epoch
is being chosen, it is no longer a test set but a second validation set, and its
numbers stop being an honest estimate.
\end{resep}

Dropout takes another route: during training every unit is switched off at random
with probability $1-p$, and the ones that survive are divided by $p$ so that the
expectation does not shift. In everyday language: at every step some members of the
team do not show up, so that no member can be relied on alone.

\begin{lemma}[Expectation and variance of the dropout mask]
\index{dropout}
Suppose $m=B/p$ with $B\sim\mathrm{Bernoulli}(p)$ and $0<p\le1$. Then
$\mathbb{E}[m]=1$ and $\operatorname{Var}(m)=(1-p)/p$. For a unit with inputs
$z_1,\dots,z_k$ and independent masks $m_1,\dots,m_k$, its output
$s=\sum_j m_jw_jz_j$ satisfies $\mathbb{E}[s]=\sum_jw_jz_j$ and
$\operatorname{Var}(s)=\frac{1-p}{p}\sum_j(w_jz_j)^2$.
\end{lemma}

\begin{proof}[Proof]
$\mathbb{E}[m]=p/p=1$ and $\operatorname{Var}(m)=\mathbb{E}[m^2]-1=p/p^2-1=(1-p)/p$.
For $s$, linearity of expectation gives the first term, and independence of the masks
gives the variance of the sum, namely $\sum_j(w_jz_j)^2\operatorname{Var}(m_j)$.
\end{proof}

That variance formula explains two things that are often forgotten. Dropout noise is
proportional to the sum of squared contributions, so units with large weights receive
relatively more noise, and dropout points in the same direction as an L2 penalty. And
that noise disappears entirely at test time, so that the model that is used differs
from the model that was trained; on very small networks that price can exceed the
benefit.

Data augmentation multiplies examples by changing the examples that already exist,
for instance shifting them or adding small noise, provided that the label does not
change. For a linear model, its relation to L2 can be computed exactly.

\begin{proposisi}[Gaussian noise on the input is the same as an L2 penalty]
\index{data augmentation}
Suppose $\hat y=\mathbf{w}^{\top}\mathbf{x}$ and $\boldsymbol{\delta}\sim
N(\mathbf{0},\sigma^2I)$ is added to the input. Then
\[
\mathbb{E}_{\boldsymbol{\delta}}\bigl[(\mathbf{w}^{\top}(\mathbf{x}+\boldsymbol{\delta})-y)^2\bigr]
=(\mathbf{w}^{\top}\mathbf{x}-y)^2+\sigma^2\|\mathbf{w}\|_2^2 .
\]
\end{proposisi}

\begin{proof}[Proof]
Expand: the linear term is $2(\mathbf{w}^{\top}\mathbf{x}-y)\,
\mathbb{E}[\mathbf{w}^{\top}\boldsymbol{\delta}]=0$, and the quadratic term is
$\mathbb{E}[(\mathbf{w}^{\top}\boldsymbol{\delta})^2]
=\mathbf{w}^{\top}\mathbb{E}[\boldsymbol{\delta}\boldsymbol{\delta}^{\top}]\mathbf{w}
=\sigma^2\|\mathbf{w}\|_2^2$.
\end{proof}

So adding small Gaussian noise to the input is the same as adding an L2 penalty with
$\lambda=\sigma^2$: the two are one idea written twice. The numerical check, for
$\mathbf{w}=(1,2)$, $\mathbf{x}=(1,0)$, $y=3$, and $\sigma=0.1$: the original
objective is $4$, the formula predicts $4+0.01\cdot5=4.05$, and a simulation of one
million noises gives $4.049107$. That difference of $0.0009$ is comparable to the
random error of the simulation, not to a mistake in the formula.

A warning for the security field: augmentation is valid only when the label is
genuinely invariant to the change that is made. Adding small noise to traffic
features is allowed, but reversing the order of packets or shifting arrival times can
turn an attack into not-an-attack. Augmentation with a wrong label is not
regularization but a lie that will keep looking good as long as the validation data
is changed in the same way.
\section{SGD, momentum, and Adam}
\tingkat{advanced}

Up to here we have treated "training" as a single box. The contents of that box
determine how fast the error falls and, on a limited step budget, how low an error is
reached. In everyday language: the right direction is not enough; what decides is how
the step is taken and how large it is.

Stochastic gradient descent replaces the full gradient by the gradient over one
minibatch of size $b$. That estimate is unbiased, and its variance shrinks
proportionally to $1/b$: $\operatorname{Var}(\hat{\mathbf{g}}_b)=\sigma_g^2/b$. So a
large minibatch reduces the wobble, but does not change the mean direction. What
changes the direction is momentum and Adam.

\begin{definisi}[Momentum and Adam]
\index{momentum}
\index{adam}
Momentum keeps a running average of the gradient:
$\mathbf{v}_t=\beta\mathbf{v}_{t-1}+\mathbf{g}_t$ and
$\mathbf{w}_t=\mathbf{w}_{t-1}-\eta\mathbf{v}_t$, with $0\le\beta<1$. Adam keeps two
moments, $\mathbf{m}_t=\beta_1\mathbf{m}_{t-1}+(1-\beta_1)\mathbf{g}_t$ and
$\mathbf{v}_t=\beta_2\mathbf{v}_{t-1}+(1-\beta_2)\mathbf{g}_t^2$ (squared per
coordinate), then uses bias correction and the update
$\mathbf{w}_t=\mathbf{w}_{t-1}-\eta\,\hat{\mathbf{m}}_t/(\sqrt{\hat{\mathbf{v}}_t}+\epsilon)$
with $\hat{\mathbf{m}}_t=\mathbf{m}_t/(1-\beta_1^t)$ and
$\hat{\mathbf{v}}_t=\mathbf{v}_t/(1-\beta_2^t)$.
\end{definisi}

That bias correction is not decoration: without it $\mathbf{m}_1=(1-\beta_1)\mathbf{g}_1$
is very small, so the first step is almost zero. Division by $\sqrt{\hat{\mathbf{v}}}$
matters even more for badly scaled problems: it makes the step size in each
coordinate roughly $\eta$, no matter whether the gradient is large or small. Momentum
only smooths the direction over time.

The numbers on the following quadratic show the difference. Take
$f(\mathbf{w})=\frac12\mathbf{w}^{\top}H\mathbf{w}$ with $H=\operatorname{diag}(1,10^4)$
and $\mathbf{w}_0=(1,0)$; the condition number of $H$ is $10^4$, which means the curve
is narrow in one direction and very flat in the other. We count steps until
$f\le10^{-8}f(\mathbf{w}_0)$.

\begin{center}
\small
\begin{tabular}{lrl}
\hline
method & rate & steps to $10^{-8}$ \\
\hline
gradient descent & $10^{-4}$ & 92099 \\
momentum ($\beta=0.9$) & $10^{-4}$ & 9132 \\
Adam & $0.1$ & 82 \\
\hline
\end{tabular}
\end{center}

Gradient descent needs $92099$ steps because its largest rate is limited by the
largest curvature, while the flat direction moves only at that rate. Momentum cuts
one digit, in line with the estimate that acceleration changes the dependence on the
condition number $\kappa$ into about $\sqrt\kappa$. Adam cuts three digits, because
per-coordinate division means the flat direction is no longer punished by the steep
one. That is not a general promise: on well-scaled problems the advantage of Adam
shrinks.

On the logistic problem of this chapter, with minibatch eight and a target training
error of $0.40$, the median of five random seeds gives SGD with rate $0.05$ a hundred
steps, whereas Adam with rate $0.1$ needs only ten. The conclusion must be careful:
the advantage is about one digit and partly comes from the larger rate, not from a
better direction.

The learning rate is usually not constant. Two common schedules are step decay,
namely the rate divided by ten at certain epochs, and cosine decay
$\eta_t=\eta_{\min}+\frac12(\eta_0-\eta_{\min})(1+\cos(\pi t/T))$: large steps at the
start so as to approach quickly, small steps at the end so as not to jump around in
the valley. With $3000$ minibatch steps and $\lambda=0.1$, a constant rate of $0.2$
gives a median training error of $0.3866$, whereas a cosine from $0.5$ to $0.01$ gives
$0.3549$.

\section{Cross validation, hyperparameter selection, and leakage}
\tingkat{advanced}

All the validation numbers in this chapter come from a set that is not used for
training. That is already correct, but it is not enough: as soon as we choose
$\lambda$ with an eye that sees the validation error, that validation error stops
being an unbiased estimate of the population error. In everyday language: if we
choose the exam questions by the grade we get, that grade no longer measures our
ability.

\kamus{cross validation}{validasi silang}{rotating the roles of the data in turn, so that
every example has once been an exam for a model that did not see it}

\begin{definisi}[$k$-fold cross validation]
\index{cross validation}
The training data is split into $k$ folds that are disjoint and of similar size. For
each $j$, the model is trained on $k-1$ folds and tested on fold $j$; the
cross-validation error is the average error over those $k$ folds, and the
hyperparameter that is chosen is the one that minimizes it. For $k=n$ this procedure
is called \emph{leave-one-out}.
\end{definisi}

Cross validation answers one weakness of a single split: at $n=60$ the validation set
contains twelve examples, and the error on twelve examples wobbles with a standard
deviation of about $0.12$. With five folds every example is tested once, but the
wobble does not disappear.

\begin{center}
\small
\begin{tabular}{rrr}
\hline
$\lambda$ & CV error (average of 5 partitions) & spread between partitions \\
\hline
0     & 0.3233 & 0.1245 \\
0.003 & 0.2400 & 0.0253 \\
0.030 & 0.2467 & 0.0217 \\
0.100 & 0.2467 & 0.0274 \\
1     & 0.2967 & 0.0298 \\
\hline
\end{tabular}
\end{center}

The first row shows two things at once: without regularization the CV error is both
the worst and the least stable, with a spread of $0.1245$, four times that of the
other rows. The second to fourth rows are almost indistinguishable: $0.2400$,
$0.2467$, and $0.2467$, with a spread of about $0.025$, and across five different
partitions the winner keeps changing. This CV curve is flat around its minimum, so
that choosing its exact value is an illusion of precision.

The consequence is real: CV selects $\lambda=0.003$ with validation error $0.2688$,
whereas the unselected $\lambda=0.1$ gives $0.2406$; that difference of $0.0282$ is
not far from one standard deviation of the CV curve. The lesson is not that cross
validation is broken, but that with sixty examples, hyperparameter selection carries
an uncertainty as large as the effect being measured. Reporting "the best $\lambda$
is $0.003$" without mentioning the width of the curve is the subtlest lie in
this chapter.

\begin{tangan}
Materials: ten cards, a pencil, paper. Write one two-digit number on each card,
shuffle, then lay out seven cards face up. Build a "rule" in the form of a list:
position one is worth this much, position two is worth that much, through position
seven.

Cover one card with your palm and guess. Because your list contains that card, your
guess is exact. Repeat for all seven positions: seven out of seven correct, and that
is your training score.

Now gather the ten cards, shuffle, lay out seven new cards, and use the same list.
Cover one card, guess, uncover. What must be visible: your list does not contain the
new cards, so your guess is pure luck. Repeat ten times with new layouts and record
how often you are exact; that is your exam score. Finally write a second rule that
uses only one property, for example "the card on the left is not larger than the card
on its right", then compare the two scores.
\end{tangan}

There is a source of error larger than that wobble and invisible in any number:
leakage, namely information from the test data entering the training process. Its
most common form is standardizing features using the mean and standard deviation of
all the data before the split is made.

The numbers, with an artificial covariate shift: if validation is shifted by $0.5$,
standardization with training statistics only gives $0.3765$, whereas pooled
statistics give $0.3361$; at a shift of $1.0$ the numbers are $0.5387$ against
$0.3819$. Pooled statistics always look better, and that is precisely the problem: they
use information about where the test data lies that is not available in the field.
Without a shift the two are equal, $0.2452$ against $0.2463$, a sign that the earlier
difference came from leakage.

The honest protocol is layered: separate train, validation, and test from the start;
compute every quantity, including means, standard deviations, decision thresholds,
and feature selection, inside the training folds only; choose hyperparameters on the
validation set or through cross validation inside the training data; touch the test
set once, then report it as it is.

The relation between capacity and the amount of data can be made quantitative. For a
finite hypothesis class, the generalization gap is controlled by the square root of
the logarithm of the number of hypotheses divided by the number of examples.

\begin{teorema}[A uniform bound for finite classes]
\index{capacity}
Suppose $\mathcal{H}$ is finite with $|\mathcal{H}|=m$, its error is 0-1, and $n$
examples are drawn independently from $\mathcal{D}$. With probability at least
$1-\delta$ over the drawing of the examples, it holds simultaneously for all
$h\in\mathcal{H}$
\[
\bigl|R(h)-\hat R_n(h)\bigr|\le\sqrt{\frac{\ln m+\ln(2/\delta)}{2n}} .
\]
\end{teorema}

\begin{proof}[Proof]
For a fixed $h$, $\hat R_n(h)$ is the average of $n$ Bernoulli variables scaled to
$[0,1]$, so Hoeffding's inequality gives
$\mathbb{P}(|R(h)-\hat R_n(h)|>t)\le2e^{-2nt^2}$. Combine over all $h$ with the union
bound: the probability that some $h$ exceeds $t$ is at most
$2m\,e^{-2nt^2}$. Setting that expression equal to $\delta$ gives
$t=\sqrt{(\ln m+\ln(2/\delta))/(2n)}$, and its complement is the statement of the
theorem.
\end{proof}

For the class that "memorizes", namely a model that gives one label to every binary
pattern over $d$ features, $m=2^d$. With $d=21$ and $\delta=0.05$ the bound is
$\sqrt{(21\ln2+\ln40)/(2n)}$, namely $0.3899$ at $n=60$ and $0.0617$ at
$n=2400$. The measured gap is $0.1338$ and $0.0097$ at the same $n$, so the bound
holds but is loose; that looseness is not a defect, because the bound must hold for
every distribution and every algorithm.

What matters is not the numbers but their shape. The gap shrinks proportionally to
$1/\sqrt n$, whereas the number of hypotheses enters through a logarithm. Quadrupling
the data only halves the gap, and increasing capacity tenfold only raises the bound by
the square root of the logarithmic increment. This is the honest reason why large
models are still useful: as long as their gap is controlled by regularization, rising
capacity does not automatically mean a rising gap.
\section{Worked example: closed-form ridge on four points}
\tingkat{intermediate}

All the formulas above can be checked by hand on a small example: four points with
one feature, $\mathbf{x}=(1,2,3,4)$ and $\mathbf{y}=(2,4,6,9)$, a model without an
intercept $\hat y=wx$, and the objective $\frac12\sum_i(y_i-wx_i)^2+\frac{\lambda}{2}w^2$.
Setting its derivative to zero gives
\[
w(\lambda)=\frac{\sum_ix_iy_i}{\sum_ix_i^2+\lambda}=\frac{64}{30+\lambda},
\]
since $\sum x_i^2=30$ and $\sum x_iy_i=2+8+18+36=64$.

\begin{center}
\small
\begin{tabular}{rrrr}
\hline
$\lambda$ & $w$ & training SSE & $|w\cdot5-10|$ \\
\hline
0   & 2.133333 & 0.466667 & 0.666667 \\
1   & 2.064516 & 0.608741 & 0.322581 \\
2   & 2.000000 & 1.000000 & 0.000000 \\
5   & 1.828571 & 3.253061 & 0.857143 \\
10  & 1.600000 & 9.000000 & 2.000000 \\
30  & 1.066667 & 34.600000 & 4.666667 \\
100 & 0.492308 & 81.255621 & 7.538462 \\
\hline
\end{tabular}
\end{center}

Figure~\ref{fig:lintasan-bobot} plots both weight trajectories against $\lambda$ on
this data: L2 shrinks without end, L1 stops exactly at zero. The last column uses one
test point at $x=5$ with true value $10$, which does not take part in determining
$w$. That is where the curve takes a U shape and its minimum falls exactly at
$\lambda=2$: there $w=2$, so that the prediction at $x=5$ is $10$ and the error is
zero. Without regularization, $w=64/30=2.1333$ is more attractive on the training data
(SSE $0.466667$ against $1.000000$) but predicts $10.667$, off by $0.667$. This is the
bias-variance trade in its barest form: the training SSE doubles, the test error
becomes zero.

\begin{figure}[htbp]
\centering
\begin{tikzpicture}
\begin{axis}[
  width=11.8cm, height=5.4cm,
  xmode=log,
  xmin=0.5, xmax=300, ymin=-0.12, ymax=2.45,
  axis lines=left,
  xtick={1,2,10,64,100},
  xticklabels={$1$,$2$,$10$,$64$,$100$},
  ytick={0,0.5,1.0,1.5,2.0},
  yticklabels={$0$,$0.5$,$1.0$,$1.5$,$2.0$},
  tick label style={font=\scriptsize},
  xlabel={regularization strength $\lambda$ (log scale)},
  ylabel={weight $w$},
  xlabel style={font=\small},
  ylabel style={font=\small},
  legend style={font=\scriptsize,draw=none,at={(0.97,0.97)},anchor=north east},
]
\addplot[domain=0.5:300,samples=200,very thick,layerindigo] {64/(30+x)};
\addlegendentry{L2: $64/(30+\lambda)$}
\addplot[domain=0.5:63.9,samples=200,very thick,layerorange] {(64-x)/30};
\addlegendentry{L1: $\max\{(64-\lambda)/30,\,0\}$}
\addplot[domain=64:300,samples=2,very thick,layerorange] {0};
\draw[densely dashed,layergrey] (axis cs:2,-0.12)--(axis cs:2,2.20);
\draw[densely dashed,layergrey] (axis cs:64,-0.12)--(axis cs:64,2.20);
\addplot[only marks,mark=*,mark size=1.8pt,layerindigo] coordinates {(2,2.0000) (64,0.6809)};
\addplot[only marks,mark=square*,mark size=1.8pt,layerorange] coordinates {(2,2.0667)};
\addplot[only marks,mark=o,mark size=3.4pt,layerorange] coordinates {(64,0)};
\node[font=\scriptsize,anchor=south] at (axis cs:2,2.24) {$\lambda=2$};
\node[font=\scriptsize,anchor=south] at (axis cs:64,2.24) {$\lambda=64$};
\node[font=\scriptsize,anchor=south west] at (axis cs:70,0.06) {L1 exactly zero};
\end{axis}
\end{tikzpicture}
\caption{Weight trajectories against regularization strength on the worked example
with four points, with $\sum x_i^2=30$ and $\sum x_iy_i=64$. The L2 penalty gives
$w=64/(30+\lambda)$, which shrinks but never touches zero. The L1 penalty gives
$w=\max\{(64-\lambda)/30,\,0\}$, which reaches zero exactly at $\lambda=64$. At
$\lambda=2$ the two are nearly the same, $2.0000$ against $2.0667$; at $\lambda=64$ L2
is still $0.6809$ whereas L1 is already $0.0000$. The dashed lines mark the two values
of $\lambda$ that are compared.}
\label{fig:lintasan-bobot}
\end{figure}

Note the right-hand wall. At $\lambda=100$, $w$ is only $0.4923$, the training SSE
explodes to $81.255621$, and the test error is $7.538462$, worse than a model that
learns nothing at all. This example also shows why searching for $\lambda$ on the
training data is misleading: the training SSE curve rises monotonically, so that it
always chooses $\lambda=0$.

\begin{tantangan}
A model reports a training error of $0.0000$ on $60$ examples and a test error of
$0.4100$. A colleague proposes three changes: (i) train longer, (ii) add features up
to degree seven, (iii) raise $\lambda$ from $0$ to $0.1$. For each proposal, predict
the sign of the change in training error and in test error, then state which one
cannot possibly help, with the reason.

Hint 1: first ask in which direction capacity moves under each proposal.

Hint 2: two of the three proposals enlarge capacity, and this model is already at zero
training error, so that enlarging capacity only widens the gap. What helps is what
narrows it, and that must raise the training error.

If you get stuck, go back to the Definition of underfitting and overfitting in the
first section.
\end{tantangan}

\kertas{Check the second column without a calculator: $64/32=2$ exactly, $64/31=2.064516$,
and $64/40=1.6$. Then compute the SSE at $\lambda=2$ yourself: the residuals are
$(0,0,0,1)$, so SSE $=1$. Finally answer: which $\lambda$ makes $w=1.28$, and what is
the prediction at $x=5$ then? Hint: solve $64/(30+\lambda)=1.28$.}
\section{Computational lab}
\tingkat{advanced}

Every number in this chapter comes from a single numpy script with no machine-learning
library. The data is deliberately small so that it can be rechecked within minutes: two
raw features, twenty-one polynomial features up to degree five, sixty training examples,
and twenty thousand validation examples from the same distribution. The labels are not
deterministic, so that there is an irreducible error of about $0.21$ which no model can
go below.

\begin{lab}
\begin{verbatim}
import numpy as np

def data(n, rng):
    x = rng.uniform(-1, 1, size=(n, 2))
    p = 1/(1 + np.exp(-4.0*(0.9 - x[:,0]**2 - x[:,1]**2)))
    return x, (rng.random(n) < p).astype(float)

def poly(x, deg=5):
    return np.column_stack([x[:,0]**i * x[:,1]**j
                            for i in range(deg+1) for j in range(deg+1-i)])

def sig(z):
    return 1/(1 + np.exp(-np.clip(z, -400, 400)))

def soft(t, l):
    return np.sign(t)*np.maximum(np.abs(t) - l, 0.0)

def er(Z, y, w, b):
    return float(np.mean((sig(Z@w + b) > 0.5) != y))

def fit(Z, y, lam, iters=30000, lr=0.05):
    n, d = Z.shape
    w = np.zeros(d); b = 0.0
    mw = np.zeros(d); vw = np.zeros(d); mb = 0.0; vb = 0.0
    for t in range(1, iters+1):
        p = sig(Z@w + b); g = (p - y)/n
        gw = Z.T@g + lam*w; gb = g.sum()
        mw = 0.9*mw + 0.1*gw; vw = 0.999*vw + 0.001*gw*gw
        mb = 0.9*mb + 0.1*gb; vb = 0.999*vb + 0.001*gb*gb
        w -= lr*(mw/(1-0.9**t))/(np.sqrt(vw/(1-0.999**t)) + 1e-8)
        b -= lr*(mb/(1-0.9**t))/(np.sqrt(vb/(1-0.999**t)) + 1e-8)
    return w, b

def fit_l1(Z, y, lam, iters=150000):
    n, d = Z.shape
    w = np.zeros(d); b = 0.0
    st = 0.5/(0.25*np.linalg.eigvalsh(Z.T@Z/n).max())
    for _ in range(iters):
        p = sig(Z@w + b); g = (p - y)/n
        w = soft(w - st*(Z.T@g + lam*w), st*lam); b -= st*g.sum()
    return w, b
\end{verbatim}
\end{lab}

The first part computes the error for eleven values of $\lambda$; the second part repeats
that with the L1 penalty through proximal gradient descent, namely a gradient step
followed by the soft threshold from the lemma above; the third part stops early; the
fourth and fifth parts compare the optimizers; the sixth to eighth parts examine cross
validation, leakage, and capacity. The actual output, cut only where it repeats:

\begin{verbatim}
d=21  n_train=60  n_validation=20000
==1 L2: lam  |w|max  train  validation
  0.0000 406.204  0.0000  0.3262
  0.0010   3.056  0.1000  0.2927
  0.0030   1.814  0.1000  0.2688
  0.0100   1.018  0.1167  0.2475
  0.0300   0.631  0.1167  0.2433
  0.1000   0.399  0.1500  0.2406
  0.3000   0.268  0.2333  0.2452
  1.0000   0.142  0.2500  0.2695
  3.0000   0.062  0.3333  0.3205
 10.0000   0.022  0.3333  0.3395
100.0000   0.003  0.3333  0.3395
==2 L1: lam  zeros/total  train  validation
  0.0020 16/21  0.1167  0.2630
  0.0050 11/21  0.1167  0.2455
  0.0100 14/21  0.1167  0.2412
  0.0200 14/21  0.1500  0.2369
  0.0500 11/21  0.1833  0.2334
  0.1000  7/21  0.2667  0.2492
  0.2000  1/21  0.3333  0.3395
==3 early stopping: epoch  train  validation  |w|max
ep=   1  0.1333  0.2481  0.15
ep=   5  0.1333  0.2429  0.38
ep=  10  0.1333  0.2432  0.48
ep=  50  0.1167  0.2454  0.77
ep= 200  0.1167  0.2541  1.38
ep= 800  0.1000  0.2879  2.85
ep=2400  0.0833  0.3109  4.81
best 0.2427 at epoch 8
==4 quadratic diag(1,1e4): steps to error 1e-8
GD=1e-4  92099 | momentum  9132 | Adam=0.1  82
==5 logistic minibatch: steps to training error 0.40 (median of 5 seeds)
sgd  lr=0.05 [100, 100, 90, 90, 100] median=100
sgd  lr=0.2  [30, 50, 40, 40, 20] median=40
adam lr=0.03 [30, 20, 20, 20, 30] median=20
adam lr=0.1  [20, 10, 20, 10, 10] median=10
==6 CV 5 partitions x 5 folds: lam  CV error (mean +- sd between partitions)
 0.000  0.3233 +- 0.1245
 0.003  0.2400 +- 0.0253
 0.030  0.2467 +- 0.0217
 0.100  0.2467 +- 0.0274
 1.000  0.2967 +- 0.0298
==7 standardization: shift  training stats  pooled stats
0.0  0.2452  0.2463
0.5  0.3765  0.3361
1.0  0.5387  0.3819
==8 capacity vs data (d=21, lam=0, nested sets)
n=   60 train=0.1667 validation=0.3005 diff=+0.1338 bound=0.3899
n=  150 train=0.2133 validation=0.2407 diff=+0.0274 bound=0.2466
n=  600 train=0.2067 validation=0.2260 diff=+0.0193 bound=0.1233
n= 2400 train=0.2058 validation=0.2155 diff=+0.0097 bound=0.0617
\end{verbatim}

Three things need to be read from that output, and all three disappoint ordinary
expectations.

First, the training error is nearly useless as a guide: it moves in the opposite
direction to the validation error. The validation curve is also not symmetric. From
$\lambda=0$ to $0.1$ its error falls by $0.086$; from $0.1$ to $10$ it rises by $0.099$
until it passes its starting point. When in doubt between two values, the smaller one
is safer.

Second, selection through cross validation is not automatically correct. CV chooses
$\lambda=0.003$, whereas validation chooses $0.1$; the difference between their CV errors
is $0.0006$, far below the spread between partitions of $0.025$, even though on the
validation data the two differ by $0.0282$.

Third, the theoretical bound is loose but its direction is right. The capacity
experiment here trains the model with a budget of sixty thousand steps on nested sets,
so that its training error is not zero. At $n=60$ the bound is $0.3899$ whereas the gap
is $0.1338$; at $n=2400$ the bound is $0.0617$ whereas the gap is $0.0097$. The ratio is
about three, so that the bound is useful for comparing settings, not for predicting
numbers.

\begin{untukdosen}
The teaching order I suggest: start from the opening numbers of the chapter, sixty
examples with training error zero and validation error $0.3262$, then ask whether that
model "learned". The answer of the class usually splits in two, and that split is the
best lecture material there is. After that the formal definitions, and let the
bias-variance decomposition be derived on the board in five lines. L1 regularization is
better introduced through the soft-threshold lemma than through a picture of the L1
ball, because students can check that lemma themselves.

The two misconceptions that appear most often. Students treat a low training error as
proof of a good model, then are surprised when its validation error is high; another
symptom is that they suggest "training longer" as the cure, when in the early-stopping
table that is precisely what does the damage. And they treat regularization as always
helpful, then install a large $\lambda$ with the conviction that it is "safer". The row
$\lambda=10$ with error $0.3395$ is the cure for both.

A blackboard activity that works: ask the class to guess, before it is computed, where
the minimum of the validation curve lies, then compute just three rows together; the
guess usually misses toward a $\lambda$ that is too small. The sign that the class is
ready to go on: they can explain why choosing $\lambda$ with the validation data makes
that number optimistic, and they can name one leakage that is invisible in the numbers.
\end{untukdosen}

\section{Summary}
\tingkat{basic}

\begin{ringkasan}
\begin{itemize}
  \item Training error and population error are computed in the same way on different
        sets; a training error of $0.0000$ with a validation error of $0.3262$ signals
        memorization.
  \item The population error decomposes into irreducible noise, squared bias, and
        variance; regularization trades bias for variance, and that trade pays only as
        long as the fall in variance is larger than the rise in bias.
  \item The L2 penalty adds $\lambda\mathbf{w}$ to the gradient, equivalent to
        multiplying the weights by $1-\eta\lambda$ at every step; L1 uses soft
        thresholding so that weights below $\lambda/H$ become exactly zero. Both have a
        dose: $\lambda=0.1$ gives $0.2406$, whereas $\lambda=10$ gives $0.3395$, worse
        than without regularization.
  \item Early stopping is regularization through the number of steps: $0.2427$ at
        epoch 8 against $0.3109$ at epoch 2400. Dropout adds noise with variance
        $(1-p)/p$ per unit, and Gaussian-noise augmentation with variance $\sigma^2$ is
        equivalent to an L2 penalty of scale $\sigma^2$.
  \item Optimizers differ on badly scaled problems: on a quadratic with condition
        number $10^4$, gradient descent needs $92099$ steps, momentum $9132$, Adam
        $82$; a cosine schedule lowers the training error from $0.3866$ to $0.3549$.
  \item Cross validation reduces the wobble without removing it: the CV curve is flat
        with a spread of $0.025$, and its choice misses by $0.0282$. Standardization
        leakage makes the numbers look better ($0.3361$ against $0.3765$) without being
        really better, and the generalization gap shrinks proportionally to $1/\sqrt n$:
        from $0.1338$ at $n=60$ to $0.0097$ at $n=2400$.
\end{itemize}
\end{ringkasan}

\section{Exercises}
\tingkat{intermediate}

\begin{latihan}
\begin{enumerate}
  \item Assess this claim: "Because the training error of our model is zero, the model
        is already good enough." True, false, or not enough evidence? Name the part that
        is mistaken and one additional number that you need.
  \item Prove the bias-variance decomposition for $y=f(x)+\varepsilon$ with
        $\operatorname{Var}(\varepsilon)=\sigma^2$: write the three terms, show that both
        cross terms vanish, and state where the independence of $D$ from $\varepsilon$
        is used.
  \item Derive $w(\lambda)=64/(30+\lambda)$ for the data in the worked example, then show
        that $w$ decreases monotonically in $\lambda$ but is never zero. Compare it with
        the L1 penalty on the same one-dimensional problem, and determine the $\lambda$
        that makes $w=0$ exactly.
  \item Compute the validation error for $c=0.95$ and $c=1.05$ in the shrinkage example
        with $\mu=2$, $\sigma=1$, $n=25$, then compare them with $c=1$. Explain why two
        deviations of equal size do not give losses of equal size.
  \item A report writes: "We chose the best $\lambda$ through cross validation, so that
        cross-validation number is an estimate of the error in the field." Fix that
        sentence, and name two additional procedures that are needed.
  \item Prove the soft-threshold lemma for $a<0$ in full, then explain why the result is
        symmetric in the sign of $a$. Hint: treat the cases $u<0$ and $u=0$ separately.
  \item Compute the number of weight-decay steps needed for a weight to become a quarter
        of its initial value if $\eta=0.1$ and $\lambda=0.02$, using the half-life
        $\ln2/(\eta\lambda)$. Check with a numpy simulation for a zero gradient.
  \item On the data of this chapter, $\lambda=0.1$ gives a training error of $0.1500$ and
        a validation error of $0.2406$. Compute the gap, compare it with the gap at
        $\lambda=0$ and $\lambda=10$, then explain why a small gap is not enough to
        conclude that a model is better.
  \item Derive the uniform bound $\sqrt{(\ln m+\ln(2/\delta))/(2n)}$ from Hoeffding's
        inequality and the union bound, then compute it for $d=10$, $n=200$, and
        $\delta=0.01$. How many times must the data be doubled for the bound to halve?
\end{enumerate}

Hints for the three hardest items. Item 3: its one-dimensional objective is
$\frac12\sum_i(y_i-wx_i)^2+\lambda|w|$; at $w=0$ its subgradient contains zero when
$\lambda\ge|\sum_ix_iy_i|=64$, so that $w=0$ exactly at $\lambda=64$. Item 6: at $u=0$
the subgradient is the interval $[-a-t,\,-a+t]$, which contains zero when $|a|\le t$.
Item 9: with $m=2^{10}$ and $\delta=0.01$ the bound is
$\sqrt{(10\ln2+\ln200)/400}=0.1749$; halving it demands four times the data.
\end{latihan}

\begin{untukpeneliti}
What is proved in this chapter: the bias-variance decomposition together with its
conditions; the equivalence of the weight-decay step with an L2 penalty; the closed form
of the minimizer of $\frac H2(w-w^\ast)^2+\lambda|w|$ and the subgradient condition that
makes L1 produce true zeros; the equivalence of Gaussian-noise augmentation with an L2
penalty for linear models; the formulas for the expectation and variance of the dropout
mask; and the uniform bound for finite hypothesis classes. What is only popular: that
regularization always helps; that Adam is always faster; that early stopping "is the
same as" L2 (it is an approximation for a quadratic objective with a small rate); that
the count of zeros under L1 measures feature importance.

Primary sources: Goodfellow, Bengio, and Courville, \emph{Deep Learning}, Chapter 5
(capacity, bias-variance), Chapter 7 (L1 and L2 penalties, augmentation, noise, early
stopping, dropout), Chapter 8 (SGD, momentum, Adam, learning-rate schedules,
hyperparameter selection). Wasserman, \emph{All of Statistics}, Chapter 13 and Chapter 20.
The uniform bound is a direct consequence of Hoeffding's inequality and the union bound.

Limits of validity: every number comes from one synthetic set with a single $60/20000$
split, so that each validation error carries an uncertainty of order $0.01$ to $0.02$.
For that reason the comparison of L1 and L2 here is not a recommendation. The uniform
bound is only for finite classes and 0-1 error.

Open questions: how large the leakage of batch normalization is, since it computes
statistics from the minibatch and so mixes information between examples; how to choose
$\lambda$ when the cost of mistakes is asymmetric; and whether a flat CV curve can be
turned into a valid interval, not merely a spread between partitions.
\end{untukpeneliti}

\peta{Chapter 26 in Goodfellow, Bengio, and Courville, \emph{Deep Learning}: Chapter 5
(capacity, overfitting and underfitting, the bias-variance decomposition), Chapter 7 (L1
and L2 penalties, weight decay, data augmentation, noise as regularization, early
stopping, dropout), Chapter 8 (SGD, momentum, Adam, learning-rate schedules,
hyperparameter selection). In Wasserman, \emph{All of Statistics}: Chapter 13 (model
selection and cross validation), Chapter 20 (bias-variance in smoothing). Backpropagation
and network training are covered in Chapter 23; honest reading of evaluation metrics,
including leakage and data splitting, continues in Chapter 25; the convex optimization
underlying penalties and subgradients is in Chapter 21.}
